

%%%%%%%%%%%%%%%
\begin{frame}{Le schéma de la base ``Messagerie''}

On part du schéma entité-association.

\vfill

\figSlideWithSize{ea-messagerie.pdf}{9cm}
\vfill


Contient tout ce qu'il faut savoir pour créer le schéma relationnel.

\end{frame}


%%%%%%%%%%%%%%%
\begin{frame}{Application de l'algorithme de normalisation}


Les dépendances données par les entités\,:
\vfill

\begin{redbullet}
  \item $idContact \to nom, pr\acute{e}nom, email$
    \item $idMessage \to contenu, dateEnvoi$
\end{redbullet}

Les dépendances données par les associations\,:
\vfill

\begin{redbullet}
        \item $idMessage \to idPr\acute{e}d\acute{e}cesseur$ (message précédent)
        \item $idMessage \to idEmetteur$ (contact)
\end{redbullet}

\vfill

Ne pas oublier la clé \texttt{(idContact, idMessage)}
de l'association \textit{à destination  de}.


\end{frame}


%%%%%%%%%%%%%%%
\begin{frame}{Le schéma de notre base}

Trois tables\,:
\vfill

\begin{redbullet}
  \item Contact (\textbf{idContact},  nom,prénon, email)
    \item Message (\textbf{idMessage},  contenu, dateEnvoi, \textit{idEmetteur},
        \textit{idPrédécesseur})
    \item Envoi (\textbf{\textit{idDestinataire}, \textit{idMessage}})
\end{redbullet}
\vfill

Clés primaires en \textbf{gras}, clés étrangères en \textit{italiques}.

\vspace*{0.5cm}

De plus, \texttt{email} est une \alert{clé secondaire} de \texttt{Contact}


\end{frame}


%%%%%%%
\begin{frame}[fragile]{Créons les tables}

Commençons par la table \texttt{Contact}

\vfill

\begin{minted}[baselinestretch=.9,
               fontsize=\small,
               framesep=2mm]{sql}
create table Contact (idContact integer not null,
                      nom varchar(30) not null,
                      prénom varchar(30)  not null,
                      email varchar(30) not null,
                      primary key (idContact),
                      unique (email)
                   );  
\end{minted}

\vfill

\alert{Noter}\,: la clause \sqlkw{unique}
sur l'attribut \texttt{email}. 

\end{frame}


%%%%%%%
\begin{frame}[fragile]{La table des messages}

\vfill

\begin{minted}[baselinestretch=.9,
               fontsize=\small,
               framesep=2mm]{sql}
create table Message (
  idMessage  integer not null,
  contenu text not null,
  dateEnvoi   datetime,
  idEmetteur int not null,
  idPrédecesseur int,
  primary key (idMessage),
  foreign key (idEmetteur) 
        references Contact(idContact),
  foreign key (idPrédecesseur) 
         references Message(idMessage)
)
\end{minted}

\vfill

\alert{Noter}\,: clauses d'intégrité référentielle,
contrainte de participation pour l'émetteur

\end{frame}


%%%%%%%
\begin{frame}[fragile]{La table des envois}

\vfill

\begin{minted}[baselinestretch=.9,
               fontsize=\small,
               framesep=2mm]{sql}
create table Envoi ( 
    idDestinataire  integer not null,
    idMessage  integer not null,
    primary key (idDestinataire, idMessage),
    foreign key (idDestinataire) 
              references Contact(idContact),
    foreign key (idMessage) 
              references Message(idMessage)
)
\end{minted}

\vfill


\alert{Noter}\,: clés typiques de la
représentation d'une association plusieurs-plusieurs.



\end{frame}
\begin{frame}{À retenir}

\'Etant donné le schéma EA, supposé correct.

\vfill

\begin{redbullet}
  \item La création des tables est \alert{directe} par application
       de l'algorithme de normalisation
   \item Certains systèmes d'aide à la  conception génèrent automatiquement ce schéma
   \item Indiquer la contrainte de clé primaire est indispensable
   \item Indiquer les contraintes de clé étrangère garantit
            l'intégrité des références
\end{redbullet}

\vfill

\end{frame}

